%%%%%%%%%%%%%%%%%%%%%%%%%%%%% Define Article %%%%%%%%%%%%%%%%%%%%%%%%%%%%%%%%%%
\documentclass{article}
%%%%%%%%%%%%%%%%%%%%%%%%%%%%%%%%%%%%%%%%%%%%%%%%%%%%%%%%%%%%%%%%%%%%%%%%%%%%%%%

%%%%%%%%%%%%%%%%%%%%%%%%%%%%% Using Packages %%%%%%%%%%%%%%%%%%%%%%%%%%%%%%%%%%
\usepackage{listings, xcolor}
\usepackage{parskip}
\usepackage{hyperref}
\usepackage{amsmath}
\usepackage{graphicx}
\usepackage[a4paper, margin = 1.5cm]{geometry}
\usepackage{float}
\usepackage{amssymb}

%%%%%%%%%%%%%%%%%%%%%%%%%%%%%%% Title & Author %%%%%%%%%%%%%%%%%%%%%%%%%%%%%%%%
\title{BDD - 21/10/25}
\author{Trentin Tommaso}
%%%%%%%%%%%%%%%%%%%%%%%%%%%%%%%%%%%%%%%%%%%%%%%%%%%%%%%%%%%%%%%%%%%%%%%%%%%%%%%

\begin{document}
  \section{Relazioni}
    \begin{itemize}
      \item (1-N): Relazioni uno-a-molti, ogni entità compare almeno 1 volta e in genrale N volte come elemento di una coppia della relazione;
      \item (N-N): Relazioni molti-a-molti; 
      \item (1-1): Relazioni uno-a-uno
    \end{itemize}
    \subsection{Relazioni Ricorsive}
      Relazioni in cui un'entità appartiene ad una relazione con un'entità dello stesso tipoo (esempio aeroporto di partenza - volo - aeroporto d'arrivo). \textbf{Chiusura transitiva} (es. tutti gli aeroporti verso cui posso volare partendo da TS con 0+ scali)
    Se ho progettato bene le cose una superchiave esiste sempre ed è l'insieme degli attributi
    \subsection{Tipo di entità debole}
      \begin{itemize}
        \item Alberghi a Udine;
        \item Vogliamo tenere traccia delle stanze libere negli alberghi;
        \item Scelgo coppia (numero, albergo) come identificativo dell'entità Stanza
        \item Stanza è un'\textit{entità debole} - Albergo entità proprietatia
        \item Relazione Stanza - Albergo è detta \textbf{Relazione identificante}
      \end{itemize}
    \subsection{Relazioni ternarie (o di grado superiore)}
      Relazioni di grado superiore al secondo sono necessarie? 
\end{document}
